% =============================================================================
\section{Ablations \& future directions}
\label{sec:ablations}

The benchmark is built around five pre-registered questions that
together probe \emph{what} drives solubility prediction quality.  Each
question was finalised before the headline numbers in
§\ref{sec:results}, and each is answered by one or more of the focused
studies that follow this section.  Rather than restate those results
here, this section serves as a reading guide: it states each question,
its hypothesis, and points to the section where it is actually
addressed.  Code for every experiment is under
\texttt{Additional\_Experiments/} in the released repository.

% -----------------------------------------------------------------------------
\subsection{Q1 -- What constitutes a good solubility method: data, representation, or model?}
\label{sec:ablations_q1}

The headline ablation disentangles three confounded levers --- training
set size, molecular representation, and model capacity --- by varying
each one while holding the others fixed.

\paragraph{Scaling with data and capacity.}
A representative method per family (LightGBM and three FastProp MLP
variants spanning $\numval{310}$\,K to $\numval{9.0}$\,M parameters) is
trained on $\{0.05, 0.10, 0.20, 0.40, 0.60, 0.80, 1.00\}$ fractions of
\scthree-train, with all other splits held fixed.  All four models
plateau on \texttt{eval}, \texttt{ood}, and \texttt{sc3\_gold} by
$N \in [\numval{12}\text{K}, \numval{25}\text{K}]$ rows, and adding
capacity does not move the asymptote.  Full curves, power-law fits and
the train-vs-test decomposition are in §\ref{sec:data_scaling}.

\paragraph{Representation ablation.}
The model is held fixed (LightGBM with the tuned \texttt{lgb\_rdkit}
hyperparameters) and the featurizer is swapped across seven choices
(RDKit, Morgan, MACCS, Mordred, Abraham, Hansen, GCN).  Per-representation
\texttt{eval}/\texttt{ood} \RMSE{}s and the resulting Tree-SHAP attribution
maps are reported in §\ref{sec:interpretability}; that section is the
authoritative source for the representation comparison and is the
right place for it because the interpretability story it sets up
(General Solubility Equation rediscovery, Abraham/LSER axes,
substructure attribution) only makes sense once the per-featurizer
\RMSE{}s are on the table.

\paragraph{Multimodal / variance-normalised loss.}
The label distribution is multimodal across solvents
(§\ref{sec:multimodal}), which motivates a per-solvent
variance-normalised MSE objective.  We have not run this comparison;
the headline metric suite already exposes the same effect through
PS-RMSE and Z-RMSE (§\ref{sec:metric_suite}), so the loss-side
experiment is deferred to follow-up work.

% -----------------------------------------------------------------------------
\subsection{Q2 -- What are the models actually learning? (Interpretability)}
\label{sec:ablations_q2}

The full interpretability study is the dedicated section
§\ref{sec:interpretability}.  It covers SHAP across the seven
representations of Q1, block-wise (solute / solvent / temperature)
attribution, top global features, per-solvent profiles, solvent
clustering, Tree-SHAP interaction values, the Abraham/LSER axis
breakdown, and graph-based atom and BRICS-substructure attribution
for the GCN.  Headline finding: LightGBM\,$+$\,RDKit independently
rediscovers the General Solubility Equation axes and Abraham/LSER
cross-terms, and the GCN learns a BRICS substructure ontology that
matches medicinal-chemistry intuition.

% -----------------------------------------------------------------------------
\subsection{Q3 -- Can other chemical properties guide solubility? (Transfer learning)}
\label{sec:ablations_q3}

Pretraining on an adjacent solvation task and fine-tuning on \scthree{}
is the test for transfer.  We use \textsc{CombiSolv-QM}
($\sim 10^6$ rows of \textsc{cosmo-rs} $\Delta G_{\text{solv}}$ at
$\numval{298.15}$~K) as the pretraining source, with pair-level
leakage filtering against every \scthree{} holdout.  The full design
--- multi-temperature and temperature-isolated setups, full and
head-only fine-tuning, fractions $\{5\%, 25\%, 100\%\}$, five seeds ---
is in §\ref{sec:transfer}.  Headline finding: QM pretraining wins on
$9{/}9$ multi-T cells and $8{/}9$ 298-K-locked cells, with the largest
data-efficiency win at the smallest fine-tuning fraction.

% -----------------------------------------------------------------------------
\subsection{Q4 -- Is lack of data the reason deep models fail?}
\label{sec:ablations_q4}

This question \emph{builds on} the scaling experiment of Q1.  Q1 sets
up the data-fraction curves; Q4 asks whether extrapolating those curves
predicts a crossing with the aleatoric floor at any reachable $N$.
We fit the saturating power law
$\mathrm{RMSE}(N) = aN^{-b} + c$ to every (model, split) curve and
solve for the training-set size $N^\star$ at which RMSE would come
within $\delta$ of $\epsA(\text{split})$.  The closed-form result and
the per-(model, split) extrapolations are in §\ref{sec:scaling_q4}.
Headline finding: $c > \epsA + 0.05$ for every (model, split)
combination, so $N^\star = \infty$ everywhere within this
representation.  The deep-model gap is representation-limited (H2),
not data-limited (H1).

% -----------------------------------------------------------------------------
\subsection{Q5 -- Can a foundation model work for solubility?}
\label{sec:ablations_q5}

Q5 ties Q2/Q3/representation together: large pretrained molecular
encoders (Uni-Mol, GROVER, SolvBERT, ChemFM, MolMerger) produce
high-dimensional learned representations, and the question is whether
any of them unlocks a scaling regime the descriptor backbones in Q1
cannot.  We have not yet run the full foundation-model sweep.  The
infrastructure for it already exists --- the transfer protocol of
§\ref{sec:transfer} provides the leakage filter and fine-tune recipe,
the scaling protocol of §\ref{sec:data_scaling} provides the curve
template, and the SHAP / GCN-attribution machinery of
§\ref{sec:interpretability} provides the post-hoc analysis --- so the
experiment reduces to swapping the trunk.  This is the most natural
next step from the conclusions of Q1 and Q4, and we flag it as the
primary outstanding ablation.
